% !TeX encoding = UTF-8
% !TeX spellcheck = en_GB
% !TeX root = ../thesis.tex

\chapter{Conclusion}
\label{chapter:conclusions}

This thesis proposed and prototyped REPUSH, a continuous session validation protocol that extends PUSH to assume a fully compromised host and that uses an LR-PUF. This chapter summarizes the main findings of this work with respect to the research questions posed in the introduction chapter and outlines directions for future work.

\paragraph{Research question 1: How must the PUSH protocol be redesigned to operate under a threat model in which the entire host, and not just the browser, is considered potentially compromised?} Chapter \ref{chapter:design} answered this by no longer treating any host-side component as trusted. The dongle manager was reduced to a translation proxy that performs no cryptographic operations and the browser also performs no cryptographic operations. Every request sent to the dongle must be signed by the server, preventing any arbitrary payload forgery. Requests are split into non-sensitive and sensitive. Non-sensitive requests are periodically re-validated, whilst sensitive requests require the user to explicitly approve them on the dongle's touch screen. Server authentication was moved away from the browser to a CA-issued certificate that the dongle can independently verify. Credentials are entered directly on the dongle's touch screen and are encrypted end-to-end so that only the server can decrypt them. These changes are what let REPUSH resist a compromised host.

\paragraph{Research question 2: How can a logically reconfigurable PUF be integrated into the session validation life-cycle, i.e. how should reconfiguration and key rotation be incorporated into the protocol's phases?} The LR-PUF integrates reconfiguration into the session life-cycle through the $map_{in}$, $map_{out}$, $reconfigure$ and $query$ operations built on top of the weak CHOICE-PUF and combined with a BCH fuzzy commitment scheme for error correction. During phase 1, $enroll_{BCH}$ is used to derive the first key pair, then during phase 2, $query_{BCH}$ is used to reconstruct the current key for signing payloads, and finally in phase 3, a reconfiguration of the LR-PUF is done before deriving the new key pair.

\paragraph{Research question 3: What security improvements does the REPUSH protocol provide over the original PUSH protocol?} REPUSH resolves two attacks that PUSH's trusted-host assumption introduced. The silent signing attack, where a compromised host could use the dongle and prove ownership thereof by using it to sign whatever the server requests without the user's knowledge, is fully prevented for sensitive requests by requesting user approval on the dongle's screen and the impact of this attack is confined to non-sensitive requests. The enrollment hijacking attack, where a compromised browser could bind a session to an attacker's own dongle, is prevented by having the dongle authenticate the server through a CA certificate and then encrypt the user credentials so that the host cannot read them. REPUSH additionally provides passive compromise detection since any unsolicited signing prompt on the dongle's screen is enough to alert the user that the host is compromised. 

\paragraph{Research question 4: What are the performance and security overheads introduced by a PUF-enabled hardware dongle (FPGA and micro-controller) in a continuous session validation protocol?} The CHOICE-PUF was measured to have a mean per-bit uniformity of 51.04\%, an overall reliability of 99.61\%, which is well within the BCH code's 4-bit correction capacity, and the BLAKE2b hashing in the LR-PUF construction was shown to improve the CHOICE-PUF's poor per-challenge uniformity. LR-PUF operations are fast. Enrollment with BCH takes on average 50.5 ms and querying takes 24.8 ms. The bulk of the observed latency comes instead from the Arduino's software implementation of Ed25519, which makes signing take around 2.8 seconds, which is slower than PUSH's 419 ms signing time. Enrollment and key rotation take around 12 seconds. None of these delays are perceptible to the user in practice, since the requests continue to be served for some seconds while a signature is pending, and enrollment or rotation are infrequent one-time operations. In short, the PUF component on the dongle is not the performance bottleneck but rather the cryptographic operations carried by the micro-controller. The price for stronger security guarantees is paid by more latency.

\section{Future Work}
Several directions remain open for future work. A large portion of REPUSH's signing latency comes from performing Ed25519 signing and verification in software on the Arduino. A dedicated cryptography chip or using a more powerful board could significantly decrease that latency. Another direction concerns using post-quantum cryptography (PQC) schemes to make the dongle more future-proof. The current threat model excludes attackers with physical access to the dongle. Future research could perhaps encrypt the SD card's content with a key derived from the CHOICE-PUF and require a PIN to decrypt the contents. Another area of improvement would be in evaluating the inter-device uniqueness of the CHOICE-PUF by comparing across multiple FPGA chips. On the protocol side, REPUSH currently leaves the distinction between sensitive and non-sensitive requests entirely up to the application owner, which risks over- or under-classification, so future work could explore adaptive or risk-scored classification, where requests could become sensitive due to some score derived from usage. To reduce friction during login, the dongle could be used as a password manager instead of having the user manually enter the username and password every time on the screen.
